\documentclass[11pt,a4paper]{article}

\usepackage[margin=2.3cm]{geometry}
\usepackage{amsmath,amssymb,amsthm}
\usepackage{mathtools}
\usepackage[hidelinks]{hyperref}
\usepackage{parskip}
\usepackage{enumitem}
\usepackage{titlesec}

\titleformat{\section}{\normalfont\large\bfseries}{\thesection}{0.7em}{}
\titleformat{\subsection}{\normalfont\normalsize\bfseries}{\thesubsection}{0.6em}{}

\newcommand{\E}{\mathbb{E}}
\newcommand{\Var}{\mathrm{Var}}
\newcommand{\Cov}{\mathrm{Cov}}
\newcommand{\WN}{\mathrm{WN}}
\newcommand{\N}{\mathcal{N}}
\newcommand{\BD}[1]{\textsc{[B\&D~§#1]}}

\title{\vspace{-2.5em}SF2943 Part A --- Method and Mathematics Companion\\
\large Electricity load forecasting for Swedish bidding zone SE\_3}
\author{}
\date{}

\begin{document}
\maketitle
\vspace{-2.5em}

\noindent\textbf{Purpose.} This is a reading aid for \texttt{project\_part\_a.ipynb}.
It collects the mathematics used in the pipeline (clean $\to$ fit $\to$ forecast)
with citations to Brockwell \& Davis (\emph{Introduction to Time Series and
Forecasting}, 2nd ed.), the course text. It does \emph{not} reproduce numerical
results or plots --- those live in the notebook.

\section{Stationarity and the decomposition model}

\subsection*{Weak stationarity \BD{1.4}}
A time series $\{X_t\}_{t \in \mathbb{Z}}$ is \emph{weakly stationary} if the
mean function is constant and the covariance depends only on the lag:
\begin{equation}
\mu_X(t) = \E[X_t] = \mu, \qquad
\gamma_X(t, t+h) = \Cov(X_t, X_{t+h}) = \gamma_X(h).
\label{eq:stationary}
\end{equation}
The \emph{autocorrelation function (ACF)} is $\rho_X(h) = \gamma_X(h)/\gamma_X(0)$.
Stationarity is what makes the ``one realization, many lags'' trick work: sample
averages along the time axis estimate population quantities. Raw electricity load
satisfies neither condition --- it has a secular trend, a strong winter/summer
cycle, a weekday/weekend cycle, and amplitude that grows with the level --- so
the cleaning stage is devoted to transforming it into something
$\hat Y_t$ for which stationarity is a defensible assumption.

\subsection*{Classical decomposition \BD{1.5}}
The course model writes
\begin{equation}
X_t = m_t + s_t + Y_t
\label{eq:decomp}
\end{equation}
where $m_t$ is a deterministic trend, $s_t$ is a deterministic seasonal component
of known period $d$ satisfying $s_{t-d} = s_t$ and
$\sum_{j=1}^d s_j = 0$ (zero-mean over one period, by convention), and $Y_t$ is
a zero-mean residual. Electricity load has two periods, so we generalise to
\begin{equation}
X_t = m_t + s_t^{(\text{week})} + s_t^{(\text{year})} + Y_t,
\qquad d_{\text{week}} = 7,\ d_{\text{year}} = 365.
\end{equation}
The cleaning procedure estimates $\hat m_t$, $\hat s_t^{(\text{week})}$,
$\hat s_t^{(\text{year})}$ and returns
$\hat Y_t = X_t - \hat m_t - \hat s_t^{(\text{week})} - \hat s_t^{(\text{year})}$.

\subsection*{Variance stabilisation: Box--Cox with $\lambda = 0$ \BD{1.5}}
The amplitude of daily load fluctuations grows with the level (winters are bigger
in absolute MW than summers). The Box--Cox transform with parameter $\lambda = 0$
is the natural log,
\begin{equation}
X_t \;\longmapsto\; \log X_t,
\end{equation}
which turns multiplicative seasonality into additive seasonality --- i.e.\ a form
that matches~\eqref{eq:decomp} exactly. Applied \emph{first}, before decomposition.

\subsection*{Harmonic regression for $\hat s_t$ \BD{1.3}}
For period $d$ we model the seasonal component as a truncated Fourier series,
\begin{equation}
s_t \;=\; a_0 + \sum_{j=1}^{k}\Bigl[a_j \cos(\lambda_j t) + b_j \sin(\lambda_j t)\Bigr],
\qquad \lambda_j \;=\; \frac{2\pi j}{d},
\label{eq:harmonic}
\end{equation}
with $k$ the number of harmonics retained. The coefficients $\{a_j, b_j\}$ are
estimated by ordinary least squares: we form the design matrix
\begin{equation}
\mathbf{Z} \;=\; \bigl[\mathbf{1},\ \cos(\lambda_1 t),\ \sin(\lambda_1 t),\ \ldots,\ \cos(\lambda_k t),\ \sin(\lambda_k t)\bigr]
\end{equation}
and compute $\hat{\boldsymbol\beta} = (\mathbf{Z}^\top \mathbf{Z})^{-1} \mathbf{Z}^\top \mathbf{y}$.
This is a linear projection onto a finite-dimensional basis of trigonometric
polynomials --- the cleanest possible setting for least squares.

\paragraph{Choices for this project.}
Yearly: $d = 365$, $k = 2$ (fundamental + first harmonic, enough to capture the
asymmetric winter peak without over-fitting).  Weekly: $d = 7$, $k = 3$
(three harmonics are enough resolution for the near-rectangular weekday/weekend
block shape). The estimated coefficients are printed in the notebook and copied
into the summary file.

\subsection*{Polynomial trend \BD{1.5}}
After removing the seasonal components we fit the remaining slowly-varying level
by a low-order polynomial:
\begin{equation}
\hat m_t \;=\; \beta_0 + \beta_1 t + \beta_2 t^2 + \cdots,
\end{equation}
also by ordinary least squares. Default degree: 1 (linear). The notebook
escalates to degree 2 (quadratic) only if the linear residuals show visible
curvature --- concretely, if the mean of any third of the residual series
exceeds 10\% of the residual standard deviation. This heuristic is a one-line
stand-in for the visual check B\&D describes.

\subsection*{Differencing: an alternative route \BD{1.5.2}}
With the backshift operator $B X_t = X_{t-1}$, define the lag-$d$ difference
\begin{equation}
\nabla_d X_t \;=\; X_t - X_{t-d} \;=\; (1 - B^d) X_t.
\end{equation}
Applied to the decomposition~\eqref{eq:decomp} with $s_t$ of period $d$,
$$\nabla_d X_t = (m_t - m_{t-d}) + (Y_t - Y_{t-d}),$$
so the seasonal component is killed exactly. A linear trend is killed by
$\nabla_1 = 1 - B$, a quadratic trend by $\nabla_1^2$. The project uses classical
decomposition as the primary route (it yields explicit $\hat m_t$ and $\hat s_t$
to report) and differencing as a sanity check: the residual ACFs from
$\hat Y_t$ and from $\nabla_1 \nabla_7 \log X_t$ should tell the same story.

\subsection*{Sample ACF and the $\pm 1.96/\sqrt{n}$ bounds \BD{1.4.1}}
The sample autocovariance is
\begin{equation}
\hat\gamma(h) \;=\; \frac{1}{n} \sum_{t=1}^{n - |h|} (x_{t+|h|} - \bar x)(x_t - \bar x),
\end{equation}
and $\hat\rho(h) = \hat\gamma(h)/\hat\gamma(0)$ is the sample ACF. The divisor
$n$ (rather than $n-h$) is used to keep the sample covariance matrix
nonnegative definite --- this matters when we later feed $\hat\gamma$ into
Yule--Walker. For an iid sequence with finite variance,
\begin{equation}
\hat\rho(h) \;\overset{\text{approx}}{\sim}\; \N\!\left(0,\ \tfrac{1}{n}\right), \qquad h \neq 0,
\end{equation}
so approximately 95\% of the $\hat\rho(h)$ lie inside $\pm 1.96/\sqrt{n}$.
These are the blue dashed lines drawn on every ACF plot in the notebook. A
correctly cleaned residual series has a sample ACF that stays mostly inside
these bounds, with no slow decay (residual trend) and no spikes at period
multiples (residual seasonality).

\subsection*{Ljung--Box portmanteau test \BD{1.6}}
The joint statistic
\begin{equation}
Q_{\mathrm{LB}} \;=\; n(n+2) \sum_{j=1}^{h} \frac{\hat\rho^2(j)}{n - j}
\end{equation}
replaces eyeballing individual $\hat\rho(h)$. Under the null that $\{X_t\}$ is
iid, $Q_{\mathrm{LB}} \overset{\text{approx}}{\sim} \chi^2(h)$. We report it at
$h = 20$ and $h = 40$. For the \emph{cleaning} residuals we \emph{expect} iid
to be rejected --- that is what justifies fitting an ARMA model in the next
stage. For the \emph{ARMA} residuals (Section~\ref{sec:arma} below) we apply
a degrees-of-freedom correction, as discussed there.

\section{ARMA identification and estimation}
\label{sec:arma}

\subsection*{ARMA($p$, $q$) and causality/invertibility \BD{3.1}}
A causal ARMA($p$, $q$) process satisfies
\begin{equation}
\phi(B) X_t \;=\; \theta(B) Z_t, \qquad \{Z_t\} \sim \WN(0, \sigma^2),
\label{eq:arma}
\end{equation}
where
$\phi(z) = 1 - \phi_1 z - \cdots - \phi_p z^p$ and
$\theta(z) = 1 + \theta_1 z + \cdots + \theta_q z^q$.
The process is \emph{causal} if every root of $\phi(z)$ satisfies $|z| > 1$
(equivalently: $X_t$ is an absolutely convergent one-sided moving average of the
innovations $\{Z_s : s \leq t\}$), and \emph{invertible} if every root of
$\theta(z)$ satisfies $|z| > 1$ (equivalently: $Z_t$ can be written as an
absolutely convergent one-sided function of $\{X_s : s \leq t\}$). Both
conditions must hold for the estimation and prediction machinery to be
well-defined; the notebook checks them on every fitted candidate by computing
the roots of $\hat\phi$ and $\hat\theta$.

\subsection*{Identification from sample ACF and PACF \BD{3.2.3, 6.2}}
The shapes of $\hat\rho$ and of the sample partial autocorrelation
$\hat\alpha$ diagnose the order:
\begin{center}
\begin{tabular}{lll}
\toprule
Process & Sample ACF & Sample PACF \\
\midrule
AR$(p)$   & tails off  & \textbf{cuts off after lag $p$} \\
MA$(q)$   & \textbf{cuts off after lag $q$} & tails off \\
ARMA$(p,q)$ & tails off & tails off \\
\bottomrule
\end{tabular}
\end{center}
``Cuts off'' means: drops inside the $\pm 1.96/\sqrt{n}$ bounds and stays there.
This is the \emph{first-pass} tool for picking $(p, q)$ candidates before any
numerical optimisation is run. The notebook uses it to select 3--5 candidate
models, which are then compared by AICC.

\subsection*{Yule--Walker estimation for AR$(p)$ \BD{5.1.1}}
Multiplying the AR equation $X_t = \phi_1 X_{t-1} + \cdots + \phi_p X_{t-p} + Z_t$
by $X_{t-j}$ and taking expectations gives, for $j \geq 1$,
$$\gamma(j) = \phi_1 \gamma(j-1) + \cdots + \phi_p \gamma(j-p),$$
and evaluating at $j = 1, \ldots, p$ yields the Yule--Walker system
\begin{equation}
\Gamma_p \boldsymbol\phi \;=\; \boldsymbol\gamma_p,
\qquad
\sigma^2 \;=\; \gamma(0) - \boldsymbol\phi^\top \boldsymbol\gamma_p,
\label{eq:yw}
\end{equation}
where
$\Gamma_p = [\gamma(i-j)]_{i,j=1}^p$,
$\boldsymbol\gamma_p = (\gamma(1), \ldots, \gamma(p))^\top$, and
$\boldsymbol\phi = (\phi_1, \ldots, \phi_p)^\top$.
The sample version replaces $\gamma$ with $\hat\gamma$:
\begin{equation}
\hat{\boldsymbol\phi} \;=\; \hat\Gamma_p^{-1} \hat{\boldsymbol\gamma}_p,
\qquad
\hat\sigma^2 \;=\; \hat\gamma(0) - \hat{\boldsymbol\phi}^\top \hat{\boldsymbol\gamma}_p.
\end{equation}
Two properties are useful: (i) $\hat\Gamma_p$ is automatically invertible
because $\hat\gamma$ is built with divisor $n$, guaranteeing nonnegative
definiteness; (ii) $\hat{\boldsymbol\phi}$ is asymptotically equivalent to MLE
for AR processes, so Yule--Walker is used in the notebook as a cross-check on
the MLE for the best AR candidate, not as the primary estimator.

\subsection*{Gaussian MLE via the innovations algorithm \BD{5.2}}
For models with $q > 0$, Yule--Walker is not enough and we use Gaussian MLE.
The trick that makes the likelihood computable is the \emph{innovations
algorithm}: letting $\hat X_j$ denote the one-step predictor of $X_j$ from
$X_1, \ldots, X_{j-1}$, and $v_{j-1} = \E[(X_j - \hat X_j)^2]$ its mean squared
error, the \emph{innovations} $U_j = X_j - \hat X_j$ are uncorrelated by
construction, so the joint Gaussian density factorises:
\begin{equation}
L(\phi, \theta, \sigma^2)
\;=\;
(2\pi)^{-n/2} \prod_{j=1}^{n} v_{j-1}^{-1/2}
\exp\!\left\{ -\frac{1}{2} \sum_{j=1}^{n} \frac{U_j^2}{v_{j-1}} \right\}.
\end{equation}
Maximising this numerically over $(\phi, \theta, \sigma^2)$ is the Gaussian
MLE; \texttt{statsmodels.tsa.arima.model.ARIMA} does exactly this under the
hood. The notebook comments every ARMA fit with a reference to this section.

\subsection*{Order selection: AICC, not AIC \BD{5.5.2}}
Among candidate models we pick the one minimising
\begin{equation}
\mathrm{AICC}
\;=\;
-2 \ln L(\hat\phi, \hat\theta, \hat\sigma^2)
+ \frac{2(p + q + 1)\,n}{n - p - q - 2},
\label{eq:aicc}
\end{equation}
where $n$ is the sample size and the ``$+1$'' on $p+q+1$ counts $\sigma^2$.
Compared to the uncorrected AIC, AICC adds the finite-sample bias term
\begin{equation}
\frac{2k(k+1)}{n - k - 1}, \qquad k = p + q + 1,
\end{equation}
which disappears as $n \to \infty$ but meaningfully penalises overparameterised
models in finite samples. This is the only information criterion the course
uses: AIC underpenalises overfitting in finite samples, and AICC is the
standard fix. The notebook computes it as
$\mathrm{AICC} = \mathrm{AIC} + 2k(k+1)/(n-k-1)$ for each candidate.

\subsection*{Ljung--Box on ARMA residuals \BD{5.3}}
After a model is fit, we recompute the portmanteau statistic on its residuals
$\hat R_t$ and compare against a $\chi^2$ distribution with \emph{corrected}
degrees of freedom:
\begin{equation}
Q_{\mathrm{LB}}(\hat R)
\;\overset{\text{approx}}{\sim}\;
\chi^2\bigl(h - p - q\bigr).
\end{equation}
The $p + q$ degrees of freedom are ``spent'' on estimating the ARMA parameters;
ignoring this correction leads to spuriously conservative $p$-values. The
notebook passes \texttt{model\_df = p + q} to
\texttt{statsmodels.stats.diagnostic.acorr\_ljungbox}. If the residual test
rejects, we either accept the failure and discuss what structure is left, or
we go back and try a different $(p, q)$ --- we do not paper over it.

\section{Linear prediction and forecasting}

\subsection*{Best linear predictor \BD{2.5}}
For a zero-mean stationary process and given $X_1, \ldots, X_n$, the best linear
(i.e.\ minimum mean squared error among linear) predictor of $X_{n+h}$ is
\begin{equation}
P_n X_{n+h} \;=\; \boldsymbol\phi_n^\top \mathbf{X}_n
\;=\; \phi_{n1} X_n + \cdots + \phi_{nn} X_1,
\end{equation}
where the coefficients satisfy the prediction equations
\begin{equation}
\Gamma_n \boldsymbol\phi_n \;=\; \boldsymbol\gamma_n(h),
\qquad
\boldsymbol\gamma_n(h) = (\gamma(h), \gamma(h+1), \ldots, \gamma(h + n - 1))^\top,
\end{equation}
with MSE
\begin{equation}
v_n(h) \;=\; \E\bigl[(X_{n+h} - P_n X_{n+h})^2\bigr]
\;=\; \gamma(0) - \boldsymbol\phi_n^\top \boldsymbol\gamma_n(h).
\end{equation}
For non-zero-mean series, predict $X_t - \mu$ and add $\mu$ back.

\subsection*{Innovations algorithm for prediction \BD{2.5.4}}
Rewriting the predictor in terms of innovations
$U_j = X_j - \hat X_j$ instead of raw observations,
\begin{equation}
\hat X_{n+1}
\;=\;
\sum_{j=1}^{n} \theta_{nj}\,\bigl(X_{n+1-j} - \hat X_{n+1-j}\bigr),
\end{equation}
yields a clean recursion because the innovations are uncorrelated by
construction --- unlike the $X_j$ themselves, which are highly correlated.
For MA and ARMA processes this is the natural route; for pure AR the
Durbin--Levinson recursion \BD{2.5.3} does the same job directly on the
observations.

\subsection*{$h$-step ARMA forecast \BD{3.3.1}}
For $n > \max(p, q)$ and $h \geq 1$,
\begin{equation}
P_n X_{n+h}
\;=\;
\sum_{i=1}^{p} \phi_i\, P_n X_{n+h-i}
\;+\;
\sum_{j=h}^{q} \theta_{n+h-1,\, j}\,\bigl(X_{n+h-j} - \hat X_{n+h-j}\bigr).
\label{eq:hstep}
\end{equation}
The second sum vanishes once $h > q$, so long-horizon forecasts follow the AR
recursion alone and decay geometrically towards the process mean. The
prediction MSE $v_n(h)$ satisfies
$$v_n(h) \;\longrightarrow\; \Var(X_t) \quad \text{as } h \to \infty,$$
reflecting the fact that the forecast carries no information about the distant
future beyond the unconditional mean. Visually, this is the characteristic
``point forecast flattens, band widens'' shape on the forecast plot in the
notebook.

\subsection*{Prediction intervals}
Under Gaussian innovations, an approximate 95\% prediction interval is
\begin{equation}
P_n X_{n+h} \pm 1.96\,\sqrt{v_n(h)}.
\end{equation}
In the notebook this interval is computed on the \emph{residual} (log-MW) scale
and then transformed to MW by adding the extrapolated seasonal and trend
components and exponentiating. Because $\exp$ is nonlinear, the MW interval is
asymmetric (log-normal) --- this is noted as an approximation but used as-is,
since the course does not introduce bias-corrected log-normal intervals.

\section{The pipeline in one page}

Bringing it all together, the notebook executes:
\begin{enumerate}[leftmargin=1.5em,itemsep=2pt]
\item $X_t$ (raw hourly MW) $\to$ daily mean (preserves MW units);
\item $\to$ $\log X_t$ (Box--Cox, $\lambda = 0$, variance stabilisation);
\item $\to$ subtract harmonic yearly seasonal $\hat s_t^{(\text{year})}$
      ($d=365$, $k=2$, OLS on~\eqref{eq:harmonic});
\item $\to$ subtract harmonic weekly seasonal $\hat s_t^{(\text{week})}$
      ($d=7$, $k=3$);
\item $\to$ subtract polynomial trend $\hat m_t$ (OLS, degree 1 or 2 by
      curvature check);
\item $\to$ residual series $\hat Y_t$ (target: stationary);
\item $\to$ diagnose via sample ACF, PACF, Ljung--Box~\BD{1.6};
\item $\to$ identify candidate ARMA$(p, q)$ from ACF/PACF shape~\BD{3.2.3};
\item $\to$ fit each candidate by Gaussian MLE~\BD{5.2}; verify causality and
      invertibility; compute AICC~\eqref{eq:aicc}; pick the minimum;
      cross-check the best AR candidate against Yule--Walker~\eqref{eq:yw};
\item $\to$ diagnose ARMA residuals via Ljung--Box with
      $\text{df} = h - p - q$~\BD{5.3};
\item $\to$ refit chosen ARMA on the training window; compute $h$-step
      forecast~\eqref{eq:hstep} for $h = 1, \ldots, 30$ and its
      $v_n(h)$-based 95\% band;
\item $\to$ invert the cleaning transforms (add trend + seasonals, then
      $\exp$) to recover MW;
\item $\to$ report RMSE, MAE, and empirical coverage of the held-out 30-day
      test window.
\end{enumerate}

\noindent Each step in the notebook is commented with the corresponding B\&D
section, so the notebook can be read as an annotated implementation of the
theory above.

\end{document}
